

\subsection{Vector Arithmetic}
\label{subsec:vector-arithmetic}

Geometric dot-product formulas in this chapter use real Euclidean vectors. For complex vectors, the Hermitian product defined below is the appropriate positive-definite inner product.

This section explains the basic arithmetic of vector and matrix computations, as well as advanced concepts such as vector/plane projection and the basis of planes (or spaces). 


\begin{tcolorbox}[title={\textbf{\tboxdef{\ref*{subsec:vector-arithmetic}} Vector Arithmetic}}]

\begin{itemize}
\item \textbf{Addition:} Given two $n \times 1$ vectors (i.e., $n$-dimensional vectors) composed of $n$ numbers each:

$\vec{a} = (a_0, a_1,\cdots, a_{n-1})$, $\vec{b} = (b_0, b_1,\cdots, b_{n-1})$

$ $

Vector addition is defined as: 

$ \vec{a} + \vec{b} = (a_0 + b_0, \text{ } a_1  +  b_1, \text{ }\cdots, \text{ } a_{n-1}  +  b_{n-1})$

$ $

\item \textbf{Dot Product:} Given two $n$-dimensional vectors:

$\vec{a} = (a_0, a_1,\ldots, a_{n-1})$, $\vec{b} = (b_0, b_1,\ldots, b_{n-1})$

$ $

Dot product is defined as: 

$\langle \vec{a}, \vec{b} \rangle =  \vec{a} \cdot \vec{b} = a_0 b_0 + a_1 b_1 +\cdots + a_{n-1} b_{n-1}$.

$\vec{a} \cdot \vec{b}$ can also be expressed as $\|\vec{a}\|\,\|\vec{b}\|\cos\theta$, where $\theta$ is the angle between $\vec{a}$ and $\vec{b}$. If $\vec{a}$ and $\vec{b}$ point in the same direction, $\vec{a}\cdot\vec{b}$ attains its maximum value ($\|\vec{a}\|\,\|\vec{b}\|$). If $\vec{a}$ and $\vec{b}$ are orthogonal, then $\vec{a}\cdot\vec{b} = 0$.

$ $


\item \textbf{Hadamard Product:} Given two $n$-dimensional vectors:

$\vec{a} = (a_0, a_1,\ldots, a_{n-1})$, $\vec{b} = (b_0, b_1,\ldots, b_{n-1})$

$ $

Hadamard product is defined as: 

$\vec{a} \odot \vec{b} = (a_0 b_0, \text{ } a_1 b_1, \text{ }\cdots, \text{ } a_{n-1} b_{n-1})$

$ $

\item \textbf{Hermitian Product:} Given two $n$-dimensional complex vectors:

$\vec{a} = (a_0 + i \cdot a'_0, \text{ } a_1+ i\cdot a'_1, \text{ }\ldots,\text{ }  a_{n-1}+ i\cdot a'_{n-1})$ 

$\vec{b} = (b_0 + i\cdot b'_0, \text{ } b_1 + i\cdot b'_1, \text{ }\ldots, \text{ } b_{n-1} + i\cdot b'_{n-1})$

$ $

Hermitian product is a dot product with the 2nd operand as a conjugate (\autoref{subsec:imaginary}): 

$\langle\langle \vec{a}, \vec{b} \rangle\rangle =  \vec{a} \cdot \overline{\vec{b}} $

$ = (a_0 + i \cdot a'_0, \text{ } a_1+ i\cdot a'_1, \text{ }\cdots,\text{ } a_{n-1}+ i\cdot a'_{n-1})  \cdot (b_0 - i\cdot b'_0, \text{ } b_1 - i\cdot b'_1, \text{ }\cdots, \text{ } b_{n-1} - i\cdot b'_{n-1})$




\end{itemize}
\end{tcolorbox}




\subsection{Various Types of Matrix}
\label{subsec:vandermonde}


\begin{tcolorbox}[title={\textbf{\tboxdef{\ref*{subsec:vandermonde}} Matrices}}]


\begin{itemize}

\item An $n \times n$ \textbf{identity matrix} and a \textbf{reverse identity matrix} are defined as:

$I_n = \begin{bmatrix}
1 & 0 & 0 & \cdots & 0\\
0 & 1 & 0 & \cdots & 0 \\
0 & 0 & 1 & \cdots & 0 \\
\vdots & \vdots & \vdots & \ddots & \vdots \\
0 & 0 & 0 & \cdots & 1 \\
\end{bmatrix}$, $I^R_n = \begin{bmatrix}
0 & \cdots & 0 & 0 & 1\\
0 & \cdots & 0 & 1 & 0 \\
0 & \cdots & 1 & 0 & 0 \\
\vdots & \iddots & \vdots & \vdots & \vdots \\
1 & 0 & 0 & \cdots & 0 \\
\end{bmatrix}$

$ $

\item The \textbf{transpose of a matrix} $X$ is defined as element-wise swapping along the diagonal line, denoted as $X^T$, which is:


$X = \begin{bmatrix}
a_1 & a_2 & a_3 & \cdots & a_n\\
b_1 & b_2 & b_3 & \cdots & b_n \\
c_1 & c_2 & c_3 & \cdots & c_n \\
\vdots & \vdots & \vdots & \ddots & \vdots \\
m_1 & m_2 & m_3 & \cdots & m_n \\
\end{bmatrix}$, $X^T = \begin{bmatrix}
a_1 & b_1 & c_1 & \cdots & m_1\\
a_2 & b_2 & c_2 & \cdots & m_2\\
a_3 & b_3 & c_3 & \cdots & m_3\\
\vdots & \vdots & \vdots & \ddots & \vdots \\
a_n & b_n & c_n & \cdots & m_n\\
\end{bmatrix}$

$ $


\item A \textbf{Vandermonde matrix} is an $(m + 1) \times (n + 1)$ matrix defined as:

$\mathit{Vander}(x_0, x_1, \cdots, x_{m}) = \begin{bmatrix}
1 & x_0 & x_0^2 & \cdots & x_0^n\\
1 & x_1 & x_1^2 & \cdots & x_1^n \\
1 & x_2 & x_2^2 & \cdots & x_2^n \\
\vdots & \vdots & \vdots & \ddots & \vdots \\
1 & x_m & x_m^2 & \cdots & x_m^{n} \\
\end{bmatrix}$

%A Vandermonde matrix is often used for representing $m$ different $n$-degree polynomials.

\end{itemize}

\end{tcolorbox}


 





\subsection{Matrix Arithmetic}
\label{subsec:matrix-arithmetic}


Matrix-to-vector multiplication and matrix-to-matrix multiplication are defined as follows:




\begin{tcolorbox}[title={\textbf{\tboxdef{\ref*{subsec:matrix-arithmetic}} Matrix Arithmetic}}]

\begin{itemize}
\item \textbf{Matrix-to-Vector Multiplication:} Given a $m \times n$ matrix $A$ and a $n$-dimensional vector $x$: 


$A = \begin{bmatrix}
 a_{1, 1} & a_{1, 2} & a_{1, 3} & \cdots & a_{1, n}\\
 a_{2, 1} & a_{2, 2} & a_{2, 3} & \cdots & a_{2, n} \\
 a_{3, 1} & a_{3, 2} & a_{3, 3} & \cdots & a_{3, n} \\
\vdots & \vdots & \vdots & \ddots & \vdots \\
 a_{m, 1} & a_{m, 2} & a_{m, 3} & \cdots & a_{m, n} \\
\end{bmatrix} = \begin{bmatrix} 
\vec{a}_{1, * } \\ 
\vec{a}_{2, * } \\ 
\vec{a}_{3, * } \\ 
\vdots\\
\vec{a}_{m, * } \\ 
\end{bmatrix}, \text{ } \vec{x} = (x_1, x_2, \cdots, x_{n})$


$ $

The result of $A \cdot \vec{x}$ is an $m$-dimensional vector computed as:

$A \cdot \vec{x} = \Big(\vec{a}_{1, * } \cdot \vec{x}, \text{ } \vec{a}_{2, * } \cdot \vec{x}, \text{ }\cdots,\text{ } \vec{a}_{m, * } \cdot \vec{x} \Big) = \left(\sum\limits_{i=1}^{n} a_{1, i} \cdot  x_i, \text{ } \sum\limits_{i=1}^{n}  a_{2, i} \cdot x_i,\cdots, \text{ } \sum\limits_{i=1}^{n} a_{m, i} \cdot x_i \right)$


$ $


\item \textbf{Matrix-to-Matrix Multiplication:} Given a $m \times n$ matrix $A$ and a $n \times k$ matrix $B$: 



$A = \begin{bmatrix}
 a_{1, 1} & a_{1, 2} & a_{1, 3} & \cdots & a_{1, n}\\
 a_{2, 1} & a_{2, 2} & a_{2, 3} & \cdots & a_{2, n} \\
 a_{3, 1} & a_{3, 2} & a_{3, 3} & \cdots & a_{3, n} \\
\vdots & \vdots & \vdots & \ddots & \vdots \\
 a_{m, 1} & a_{m, 2} & a_{m, 3} & \cdots & a_{m, n} \\
\end{bmatrix}$, $B = \begin{bmatrix}
 b_{1, 1} & b_{1, 2} & b_{1, 3} & \cdots & b_{1, k}\\
 b_{2, 1} & b_{2, 2} & b_{2, 3} & \cdots & b_{2, k} \\
 b_{3, 1} & b_{3, 2} & b_{3, 3} & \cdots & b_{3, k} \\
\vdots & \vdots & \vdots & \ddots & \vdots \\
 b_{n, 1} & b_{n, 2} & b_{n, 3} & \cdots & b_{n, k} \\
\end{bmatrix}$

$ $

The result of $A \cdot B$ is a $m \times k$ matrix computed as:

$A \cdot B = \begin{bmatrix}
\sum\limits_{i=1}^{n}a_{1, i}b_{i, 1} & \sum\limits_{i=1}^{n}a_{1, i}b_{i, 2} & \sum\limits_{i=1}^{n}a_{1, i}b_{i, 3} & \cdots & \sum\limits_{i=1}^{n}a_{1, i}b_{i, k} \\
\sum\limits_{i=1}^{n}a_{2, i}b_{i, 1} & \sum\limits_{i=1}^{n}a_{2, i}b_{i, 2} & \sum\limits_{i=1}^{n}a_{2, i}b_{i, 3} & \cdots & \sum\limits_{i=1}^{n}a_{2, i}b_{i, k} \\
\sum\limits_{i=1}^{n}a_{3, i}b_{i, 1} & \sum\limits_{i=1}^{n}a_{3, i}b_{i, 2} & \sum\limits_{i=1}^{n}a_{3, i}b_{i, 3} & \cdots & \sum\limits_{i=1}^{n}a_{3, i}b_{i, k} \\
\vdots & \vdots & \vdots & \ddots & \vdots \\
\sum\limits_{i=1}^{n}a_{m, i}b_{i, 1}
& \sum\limits_{i=1}^{n}a_{m, i}b_{i, 2} & \sum\limits_{i=1}^{n}a_{m, i}b_{i, 3} & \cdots & \sum\limits_{i=1}^{n}a_{m, i}b_{i, k} \\
\end{bmatrix}$


\end{itemize}
\end{tcolorbox}

Given the above definitions of matrix and vector arithmetic, the following algebraic properties can be derived:

\begin{tcolorbox}[title={\textbf{\tboxtheorem{\ref*{subsec:matrix-arithmetic}} Matrix Arithmetic Properties}}]

\begin{itemize}

\item \textbf{Associative:} 

$(AB)C = A(BC)$

$A(Bx) = (AB)x$


$ $

 
\item \textbf{Distributive:} 

$A(x + y) = Ax + Ay$

%$A(x \odot y) = Ax \odot Ay $

%$A\langle x, y\rangle = \langle Ax, Ay\rangle$

%$A\langle \langle x, y\rangle\rangle = \langle\langle Ax, Ay\rangle\rangle$

$ $

These linearity identities do not assert preservation of dot, Hermitian, or Hadamard products. For example, with $A=2I$, $x=y\ne0$, one generally has $A(x\odot y)\ne(Ax)\odot(Ay)$. This is a failure of a universal identity, not an inequality valid for every input: $A=I$ does preserve the Hadamard product. Expressions such as $A\langle x,y\rangle$ are scalar multiples of a matrix, whereas $\langle Ax,Ay\rangle$ is a scalar; they must not be compared as two outputs of the same vector map.

$ $

\item \textbf{NOT Commutative:} 

Matrix multiplication is not commutative in general. Both products must first be defined and have compatible shapes; some matrices commute, for example $AI=IA$. For a column vector $x$, the expression $xA$ may be undefined even when $Ax$ is defined.


\end{itemize}
\end{tcolorbox}

\begin{proof}

$ $

The properties described in \ref*{subsec:matrix-arithmetic} can be demonstrated by expanding the formulas on both sides of each equation using a variable representation for each element in the vectors/matrices and comparing the resulting formulas. We leave this expansion as an exercise for the reader.

\end{proof}


\subsection{Projection}
\label{subsec:projection}

In this subsection the scalar field is $\mathbb{R}$. Every vector onto which we project, and every vector in an orthogonal basis, is assumed nonzero.

We consider orthogonal projection onto a line (vector projection) and onto a higher-dimensional subspace.



\begin{figure}[!h]
    \centering
    \subfloat[Vector Projection \href{https://upload.wikimedia.org/wikipedia/commons/thumb/9/98/Projection_and_rejection.png/200px-Projection_and_rejection.png}{(Source)}]{
  \includegraphics[width=0.25\linewidth]{figures/vector-projection.png}

  \label{fig:vector-projection}
}
    \subfloat[Orthogonal (Plane/Subspace) Projection \href{https://www.researchgate.net/publication/333616378/figure/fig2/AS:766235117617152@1559696096437/Geometric-illustration-of-the-orthogonal-projection-operator-P-A-vector-x-R-m-is.png}{(Source)}]{
  \includegraphics[width=0.4\linewidth]{figures/orthogonal-projection.png}

  \label{fig:orthogonal-projection}
}
  \caption{}
\end{figure}


\para{Vector Projection:} For real vectors $\vec a$ and $\vec b\ne0$, the signed scalar component along $\vec b$ is $c=(\vec a\cdot\vec b)/\|\vec b\|$. Thus
\[
\textsf{Proj}_{\vec b}(\vec a)=c\frac{\vec b}{\|\vec b\|}=\frac{\vec a\cdot\vec b}{\|\vec b\|^2}\vec b.
\]
For $\vec a\ne0$, the length of this projection is $\|\vec a\|\,|\cos\theta|$; $\|\vec a\|\cos\theta$ is the signed component, which can be negative. Figure~\ref{fig:vector-projection} illustrates the projection.

\para{Orthogonal Projection:} Given the vector $\vec{x}$ and a set of mutually orthogonal vectors $\vec{p}_0, \vec{p}_1, \cdots, \vec{p}_{n-1}$ that span the subspace $P$, the orthogonal projection $\textsf{Proj}_P(\vec{x})$ measures how much of $\vec{x}$ lies in $P$ (i.e., the component of $\vec{x}$ within $P$). In the example, $\vec{x}$’s projection onto $P$ (red arrow) equals the sum of the projections of $\vec{x}$ onto each orthogonal basis vector $\vec{p}_i$ spanning $P$: $\textsf{Proj}_P(\vec{x}) = \sum\limits_{i=0}^{n-1}\textsf{Proj}_{\vec{p}_i}(\vec{x})$. This computation can be viewed as expressing $\vec{x}$ in a coordinate system defined by the $n$ orthogonal vectors.

\begin{tcolorbox}[title={\textbf{\tboxdef{\ref*{subsec:projection}} Vector and Orthogonal Projections}}]

\begin{itemize}

\item \textbf{Vector Projection:} Given two vectors $\vec{a}$ and $\vec{b}$ in the same vector space, the vector projection of $\vec{a}$ on $\vec{b}$ is:

$\textsf{Proj}_{\vec{b}}(\vec{a}) = \vec{a}_p  = \dfrac{\vec{a} \cdot \vec{b}}{||b||^2} \vec{b}$ \textcolor{red}{ $\rhd$ where $||\vec{a}_p|| = ||\vec{a}||\,|\cos \theta|$}

$ $

\item \textbf{Orthogonal Basis:} If the $n$-dimensional plane (or subspace) $P$ is spanned by the mutually orthogonal $n$-dimensional vectors $\vec{p}_0, \vec{p}_1, \cdots, \vec{p}_{n-1}$, 

then the matrix $P = \begin{bmatrix} 
\vec{p}_0\\
\vec{p}_1\\
\vdots\\
\vec{p}_{n-1}\\
\end{bmatrix}
$
is defined to be an orthogonal basis of plane $P$. 


\item \textbf{Orthogonal Projection:} Given the orthogonal basis matrix $P = \begin{bmatrix} 
\vec{p}_0\\
\vec{p}_1\\
\vdots\\
\vec{p}_{n-1}\\
\end{bmatrix}
$, 

vector $\vec{a}$'s orthogonal projection on $P$ is:

$\textsf{Proj}_P(\vec{a}) = \sum\limits_{i=0}^{n-1} \textsf{Proj}_{\vec{p}_i}(\vec{a})$



\end{itemize}


\end{tcolorbox}


Based on the definition of orthogonal projection, the following properties are derived:  

\para{Orthogonal Basis:} In an $n$-dimensional vector space, any mutually orthogonal $n$ vectors in the vector space span the subspace $P$ that is identical to the entire vector space. Further, the orthogonal projection of any vector in the vector space on $P$ is guaranteed to be a unique vector.

\para{Non-orthogonal Basis:} Let the rows of the invertible matrix $P$ be a basis $\vec p_0,\ldots,\vec p_{n-1}$ of $\mathbb R^n$. If $\vec v$ is a column vector, its coordinates in this basis are $(P^T)^{-1}\vec v$, since $\vec v=P^T\vec c$; in general $P\vec v$ is a different invertible linear transformation, not these basis coordinates. Orthogonal projection onto the span remains well-defined even when the chosen basis is non-orthogonal. For a full-column-rank real matrix $Q$ whose columns span a subspace, it is $Q(Q^TQ)^{-1}Q^T\vec v$. For a basis of the whole space the projection is simply $\vec v$. Only the sum of individual vector-projection formula requires an orthogonal basis.

\begin{tcolorbox}[title={\textbf{\tboxtheorem{\ref*{subsec:projection}} Uniqueness of Transformed Vectors}}]

\begin{itemize}
\item  \textbf{Orthogonal Basis:} If some $n$ vectors are a orthogonal basis of the plane $P$ in the $n$-dimensional vector space, then $P$ is the same as the entire vector space, and $\textsf{Proj}_P(v)$ for every vector $\vec{v}$ in the vector space is guaranteed to be a unique vector.
\item \textbf{Non-orthogonal Basis:} For a basis stored as the rows of $P$, every vector $\vec v$ has the unique coordinate vector $(P^T)^{-1}\vec v$. The orthogonal projection onto the entire space spanned by this basis is $\vec v$ itself.
\end{itemize}
\end{tcolorbox}




\subsection{Basis of a Polynomial Ring}
\label{subsec:polynomial-ring-basis}

The ring $\mathbb Z[X]/(X^n+1)$ is a free $\mathbb Z$-module of rank $n$, not a vector space over $\mathbb Z$. A basis is a list $f_0,\ldots,f_{n-1}$ for which each element has a unique integer linear combination. Equivalently, require:

\begin{itemize}

\item \textbf{Linear Independence}: If $\sum_i z_i f_i=0$ with $z_i\in\mathbb Z$, then every $z_i=0$.
\item \textbf{Spanning the Polynomial Ring:} A linear combination of the polynomials in the basis set can express any polynomial in the polynomial ring
\end{itemize}

Note that for a $(n-1)$-degree polynomial ring, the number of polynomials that form a basis of the polynomial ring is exactly $n$.


\subsection{Isomorphism between Polynomials and Vectors over Integers (and Ring)}
\label{subsec:poly-vector-transformation}



Let $n$ be a power of $2$, let $t$ be prime with $2n\mid t-1$, and fix the $n$ distinct roots $x_0,\ldots,x_{n-1}$ of $X^n+1$ in $\mathbb Z_t$. Define $\sigma$ by evaluation at these roots. Distinctness gives interpolation of degree less than $n$ (\autoref{subsec:polynomial-interpolation-field-condition}); being roots makes evaluation well-defined on the quotient and compatible with multiplication. Evaluation at arbitrary distinct nodes would give a linear bijection of coefficient representatives, but not generally the asserted ring isomorphism. Explicitly:

$\sigma: f(X) \in \mathbb{Z}_t[X] / (X^n + 1) \text{ } \longrightarrow \text{ } (f(x_0), f(x_1), f(x_2), \cdots, f(x_{n-1})) \in \mathbb{Z}^n_t$

$ $

Now, we will explain why the mapping $\sigma$ is isomorphic, which means that $\sigma$ is a bijective one-to-one mapping from $\mathbb{Z}_t[X] / (X^n + 1)$ to $\mathbb{Z}^n_t$, and it preserves the algebraic operations $(+, \cdot)$ (i.e., $\sigma$ is a homomorphism for addition and multiplication).

$ $

\para{Bijective:} In the $(n-1)$-degree polynomial ring, a list of $y$ values evaluated at some statically chosen $n$ distinct $x$ coordinates defines a unique polynomial because, algebraically, there exists only one $(n-1)$-degree (or a lesser degree) polynomial that passes through each given set of $n$ distinct $(x, y)$ coordinates. We proved this in Lagrange Polynomial Interpolation (Theorem~\ref*{sec:polynomial-interpolation} in \autoref{sec:polynomial-interpolation}). 

$ $

\para{Homomorphic:} The homomorphism of the mapping $\sigma$ on the $(+, \cdot)$ operations means that the following two relationships hold: 

$\sigma(f_a(X) + f_b(X)) = \sigma(f_a(X)) + \sigma(f_b(X))$

$\sigma(f_a(X) \cdot f_b(X)) = \sigma(f_a(X)) \odot \sigma(f_b(X))$ \textcolor{red}{\text{ } \# $\odot$ is Hadamard vector multiplication (Summary~\ref*{subsec:vector-arithmetic})}

$ $

To prove our $\sigma$ mapping's homomorphism, let's denote the input polynomials $f_a(X)$, $f_b(X)$, and their $\sigma$-mapped output vectors as follows:


$f_a(X) = a_0 + a_1X + a_2X^2 + \cdots + a_{n-1}X^{n-1}$

$ \sigma(f_a(X)) = \left({f_a(x_0)}, {f_a(x_1)}, {f_a(x_2)}, \cdots, {f_a(x_{n-1}}) \right)  = \left(\sum\limits_{i=0}^{n-1}a_i(x_0)^i, \sum\limits_{i=0}^{n-1}a_i(x_1)^i, \sum\limits_{i=0}^{n-1}a_i(x_2)^i, \cdots, \sum\limits_{i=0}^{n-1}a_i(x_{n-1})^i\right)$


$ $

$f_b(X) = b_0 + b_1X + b_2X^2 + \cdots + b_{n-1}X^{n-1}$

$ \sigma(f_b(X)) = \left({f_b(x_0)}, {f_b(x_1)}, {f_b(x_2)}, \cdots, {f_b(x_{n-1}}) \right) = \left(\sum\limits_{i=0}^{n-1}b_i(x_0)^i, \sum\limits_{i=0}^{n-1}b_i(x_1)^i, \sum\limits_{i=0}^{n-1}b_i(x_2)^i, \cdots, \sum\limits_{i=0}^{n-1}b_i(x_{n-1})^i\right)$

$ $

Given the above setup, we can see that $\sigma$ preserves homomorphism on the $(+)$ operation as follows:

$\bm{\sigma}\bm{(}f_a(X) + f_b(X)\bm{)} = \bm{\sigma\bigg(}(a_0 + b_0) + (a_1 + b_1)X + (a_2 + b_2)X^2 + \cdots + (a_{n-1} + b_{n-1})X^{n-1}\bm{\bigg)}$

$= \left(\sum\limits_{i=0}^{n-1}(a_i + b_i)(x_0)^i, \text{ }\sum\limits_{i=0}^{n-1}(a_i + b_i)(x_1)^i, \text{ }\sum\limits_{i=0}^{n-1}(a_i + b_i)(x_2)^i, \cdots, \text{ }\sum\limits_{i=0}^{n-1}(a_i + b_i)(x_{n-1})^i\right)$

$= \left(\sum\limits_{i=0}^{n-1}a_i(x_0)^i, \text{ }\sum\limits_{i=0}^{n-1}a_i(x_1)^i, \text{ }\sum\limits_{i=0}^{n-1}a_i(x_2)^i, \cdots, \text{ } \sum\limits_{i=0}^{n-1}a_i(x_{n-1})^i\right)$

\text{ } $+ \left(\sum\limits_{i=0}^{n-1}b_i(x_0)^i, \text{ } \sum\limits_{i=0}^{n-1}b_i(x_1)^i, \text{ } \sum\limits_{i=0}^{n-1}b_i(x_2)^i, \cdots, \text{ } \sum\limits_{i=0}^{n-1}b_i(x_{n-1})^i\right)$

$= \bm{\sigma(}f_a(X)\bm{)} + \bm{\sigma(}f_b(X)\bm{)}$

$ $

Also, we can see that $\sigma$ preserves homomorphism on the $(\cdot)$ operation as follows:


$\bm{\sigma(}f_a(X)\cdot f_b(X)\bm{)} = \bm{\sigma\bigg(}\left(\sum\limits_{i=0}^{n-1}a_iX^i\right)\cdot\left( \sum\limits_{i=0}^{n-1}b_iX^i\right)\bm{\bigg)}$

$ $

$= \Bigg( \left(\sum\limits_{i=0}^{n-1}a_ix_0^i\right)\left( \sum\limits_{i=0}^{n-1}b_ix_0^i\right), \text{ }
\left(\sum\limits_{i=0}^{n-1}a_ix_1^i\right)\left( \sum\limits_{i=0}^{n-1}b_ix_1^i\right), \text{ }
\left(\sum\limits_{i=0}^{n-1}a_ix_2^i\right)\left( \sum\limits_{i=0}^{n-1}b_ix_2^i\right),$

$  \text{ } \text{ } \text{ } \text{ } \text{ } \text{ }
\cdots,
\left(\sum\limits_{i=0}^{n-1}a_ix_{n-1}^i\right)\left( \sum\limits_{i=0}^{n-1}b_ix_{n-1}^i\right)
\Bigg)$

$ $

$ = \Bigg(\sum\limits_{i=0}^{n-1}a_i(x_0)^i, \text{ } \sum\limits_{i=0}^{n-1}a_i(x_1)^i, \text{ } \cdots, \text{ } \sum\limits_{i=0}^{n-1}a_i(x_{n-1})^i\Bigg)$

$  \text{ } \text{ } \text{ } \text{ }  \odot  \Bigg(\sum\limits_{i=0}^{n-1}b_i(x_0)^i, \text{ } 
\sum\limits_{i=0}^{n-1}b_i(x_1)^i, \text{ }
\cdots, \text{ } \sum\limits_{i=0}^{n-1}b_i(x_{n-1})^i\Bigg)$

$ $

$= \bm{\sigma(}f_a(X)\bm{)} \odot \bm{\sigma(}f_b(X)\bm{)}$

$ $

In summary, $\sigma$ preserves the following homomorphism:

$\sigma(f_a(X) + f_b(X)) = \sigma(f_a(X)) + \sigma(f_b(X))$

$\sigma(f_a(X) \cdot f_b(X)) = \sigma(f_a(X)) \odot \sigma(f_b(X))$

$ $

However, for $\sigma(f_a(X) \cdot f_b(X)) = \sigma(f_a(X)) \odot \sigma(f_b(X))$, we need further reasoning to justify that this relation holds in polynomial rings, which is explained below.


\para{Polynomial Ring Reduction:} Suppose that we did not have the polynomial ring setup $X^n + 1$. Then, if we multiply $f_a(X)$ and $f_b(X)$, then $f_a(X) \cdot f_b(X)$ may become a new polynomial whose degree is higher than $n-1$. This higher-degree polynomial would still decode into the expected correct vector. Suppose the following:

$\sigma(f_a(X)) = (f_a(x_0), f_a(x_1), \cdots, f_a(x_{n-1})) = (v_0, v_1, \cdots, v_{n-1})$

$\sigma(f_b(X)) = (f_b(x_0), f_b(x_1), \cdots, f_b(x_{n-1})) = (u_0, u_1, \cdots, u_{n-1})$

$ $

Then, the following is true: 

$\sigma(f_a(X) \cdot f_b(X)) = (f_a(x_0)\cdot f_b(x_0), f_a(x_1)\cdot f_b(x_1), \cdots, f_a(x_{n-1})\cdot f_b(x_{n-1})) = (v_0u_0, v_1u_1, \cdots, v_{n-1}u_{n-1})$

$= (v_0, v_1, \cdots, v_{n-1}) \odot (u_0, u_1, \cdots, u_{n-1}) $

$ $

As shown above, even if $f_a(X) \cdot f_b(X)$ results in a polynomial with a degree higher than $n-1$, it can be decoded into the expected correct vector. However, the $\sigma$ mapping loses the property of isomorphism between a polynomial and a vector because if a polynomial's degree is higher than $n-1$, then there can be more than 1 polynomial that passes through the given $n$ distinct $X$ coordinates: $\{x_0, x_1, \cdots, x_{n-1}\}$. This is a problem because, if the $\sigma$ mapping supports only polynomial-to-vector mappings and not vector-to-polynomial mappings, then we cannot convert vectors into polynomials in the first place and do isomorphic computations. Another minor issue is that if the polynomial degree term is higher than $n-1$, then the computational overhead of decoding (i.e., polynomial evaluation) becomes larger than before. 

To resolve these two minor issues, we let the $n$ distinct $X$ coordinates of evaluation be the solutions of the polynomial ring modulo $X^n + 1$ (where $n$ is some power of 2) and reduce $f_a(X) \cdot f_b(X)$ to a new polynomial modulo $X^n + 1$ whose degree is at most $n - 1$. Let $f_{ab}(X) = f_a(X) \cdot f_b(X)$, and $f'_{ab}(X)$ be the reduced polynomial such that $f_{ab}(X) = Q(X)\cdot (X^n + 1) + f'_{ab}(X)$ for some quotient polynomial $Q(X)$. Then, as illustrated in Summary \ref*{subsubsec:polynomial-ring-discuss} (\autoref{subsubsec:polynomial-ring-discuss}), $f_{ab}(X)$ and $f'_{ab}(X)$ evaluate to the same value if they are evaluated at the roots of $X^n + 1$ (by zeroing out the $Q(X)$ term). Therefore, if we let the $n$ distinct evaluating points $\{x_0, x_1, \cdots, x_{n-1}\}$ be the roots of $X^n + 1$, then we can ensure that the decoded vector of $f'_{ab}(X)$ is identical to that of $f_{ab}(X)$, which we expect. Therefore, we can replace the higher-degree polynomial $f_{ab}(X)$ with the reduced polynomial $f'_{ab}(X)$ and continue with any further polynomial additions or multiplications using $f'_{ab}(X)$. Also, by applying polynomial ring reduction, we can enhance the computational efficiency of polynomial addition and multiplication, as well as preserve the isomorphism of the $\sigma$ mapping. Therefore, we can freely convert between vectors \& polynomials and perform additions and multiplications. 

For applying this polynomial ring reduction, the polynomial modulus can be a degree-$n$ polynomial that splits into $n$ distinct roots in the coefficient field. In practice, we often choose $X^n + 1$ as the polynomial ring modulus, which is the $(\mu=2n)$-th cyclotomic polynomial (\autoref{subsec:cyclotomic-def}). The reason we let the polynomial ring modulus be a cyclotomic polynomial (especially the $(\mu=2n)$-th cyclotomic polynomial, $X^n + 1$) is that its $n$ distinct roots are well-defined (i.e., primitive $(\mu=2n)$-th roots of unity) and thus can be quickly identified even when $n$ is large.

$ $

\para{Polynomial Coefficient Modulo Reduction:} In addition, we often reduce the polynomial coefficients based on some modulus $t$ to keep the size of the coefficients lower than a certain limit for the purpose of computational efficiency. Suppose two polynomials $f_c(X)$ and $f_d(X)$ have congruent coefficients modulo $t$ as follows:

$f_c(X) = w_0 + w_1\cdot X + w_2\cdot X^2 + \cdots + w_{n-1}\cdot X^{n-1}$

$f_d(X) = w'_0 + w'_1\cdot X + w'_2\cdot X^2 + \cdots + w'_{n-1}\cdot X^{n-1} $

$w_i \equiv w'_i \bmod t$

$ $

Then, their evaluated value $f_c(x)$ and $f_d(x)$ for any $x \in \mathbb{Z}$ is guaranteed to be congruent modulo $t$, as shown below:

$f_c(x) = w_0 + w_1\cdot x + w_2\cdot x^2 + \cdots + w_{n-1}\cdot x^{n-1}$

$\equiv (w_0 + w_1\cdot x + w_2\cdot x^2 + \cdots + w_{n-1}\cdot x^{n-1}) \bmod t $

$\equiv (w_0  \bmod t) + (w_1  \bmod t)\cdot x + (w_2  \bmod t)\cdot x^2 + \cdots + (w_{n-1} \bmod t)\cdot x^{n-1} $

$\equiv w'_0 + w'_1\cdot x + w'_2\cdot x^2 + \cdots + w'_{n-1}\cdot x^{n-1} $

$= f_d(x)$

\para{Summary:} Since $\sigma$ is bijective and homomorphic, $\sigma$ is an isomorphic mapping between the $(n-1)$-degree polynomial ring $\mathbb{Z}_t[X]/(X^n + 1)$ and the $n$-dimensional vector space $\mathbb{Z}_t^n$. 


$ $


\subsubsection{Finding Appropriate Modulus $t$} 
\label{subsubsec:poly-vector-transformation-modulus}

To isomorphically evaluate a polynomial in $\mathbb{Z}_t[X]/(X^n + 1)$ into an $n$-dimensional vector, we need to evaluate the polynomial at $n$ distinct roots of $X^n + 1 \bmod t$. However, $X^n + 1 \bmod t$ does not have $n$ distinct roots for all combinations of (degree, modulus) $ = (n, t)$. For example, if $n = 2$ and $t = 3$, then $X^2+ 1 \not\equiv 0 \bmod 3$ for any possible values of $X = \{0, 1, 2\}$. Therefore, our goal is to find a satisfactory $t$ given a fixed $n$ such that $n$ distinct roots of $X^n + 1 \pmod t$ exist in order to use the isomorphic $\sigma$ mapping. 

We start with two constraints: (1) choose $t$ to be a prime number; (2) ensure $t-1$ is a multiple of $2n$.

We learned from Fermat's Little Theorem in Theorem~\ref*{subsec:order-theorem}.4 (\autoref{subsec:order-theorem}) the following: $a^{t - 1} \equiv 1 \bmod t$ if and only if $a$ and $t$ are co-prime. This means that if $t$ is a prime, then $a^{t - 1} \equiv 1 \bmod t$ for all $a \in \mathbb{Z}_t^{\times}$ (i.e., $\mathbb{Z}_t$ without $\{0\}$). Suppose $g$ is the generator of $\mathbb{Z}_t^{\times}$ whose powered values generate all elements of $\mathbb{Z}_t^{\times}$. Then, $\textsf{Ord}_{\mathbb{Z}_t}(g) = t - 1$ and $g^{t - 1} \equiv 1 \bmod t$. Since $t - 1 = k \cdot 2n$ for some $k$, $g^{k\cdot2n} \equiv (g^{k})^{2n} \equiv 1 \bmod t$. Then, $\textsf{Ord}_{\mathbb{Z}_t}(g^{k}) \leq 2n$. However, since $\textsf{Ord}_{\mathbb{Z}_t}(g) = t - 1$, for all integers $a$ such that $1\le a < t - 1 = k \cdot 2n$, $g^a \not\equiv 1 \bmod t$. In other words, for all integers $b$ such that $1\le b < 2n$, $(g^k)^b \not\equiv 1 \bmod t$. Thus, $\textsf{Ord}_{\mathbb{Z}_t}(g^{k}) = 2n$.

Let $c = g^k$. Since $\textsf{Ord}_{\mathbb{Z}_t}(c) = 2n$, $c^{2n} \equiv 1 \bmod t$. In other words, $(c^n)^2 \equiv 1 \bmod t$. 
Now, $c^n$ can be only 1 or -1. The reason is that in the relation $X^2 \equiv 1 \bmod t$, $X$ can be mathematically only $1$ or $-1 \equiv t - 1 \bmod t$. If we substitute $X = c^n$, then $c^n$ can be only $1$ or $-1 \equiv t - 1 \bmod t$. But $\textsf{Ord}_{\mathbb{Z}_t}(c) = 2n$, thus $c^n$ cannot be 1 (because $c^n=1$ would contradict $\operatorname{ord}(c)=2n$). Thus, $c^n$ can be only $-1 \equiv t - 1 \bmod t$. If $c^n = -1 \equiv t - 1 \bmod t$, then $c$ is the root of $X^n + 1 \bmod t$, because $X^n + 1 = c^n + 1 \equiv (t - 1) + 1 \equiv 0 \bmod t$.

In conclusion, given a cyclotomic polynomial $X^n + 1$, if we choose a prime $t$ such that $t - 1 = k \cdot 2n$ for some integer $k$, then one root of $X^n + 1 \bmod t$ is: $X = c = g^k$. 

Once we have found one root of $X^n + 1$, then we can derive all $n$ distinct roots of $X^n + 1$. Suppose $\omega$ is one root of $X^n + 1 \bmod t$. Then we derive the following:

$\omega^n + 1 \equiv 0 \bmod t$

$\omega^n \equiv t - 1 \bmod t$

$\omega^{2n} \equiv (t - 1)\cdot(t - 1) \equiv  t^2 - 2t + 1 \equiv   1 \bmod t$

Here $\omega^{2n}=1$ and $\omega^n=-1\ne1$ in $\mathbb Z_t$, with odd prime $t$. Since $n$ is a power of two, every proper divisor of $2n$ divides $n$. Hence the order of $\omega$ cannot be a proper divisor, so $\operatorname{ord}(\omega)=2n$. This inference uses the power-of-two assumption; for general $n$, a root of $X^n+1$ need not be primitive.

Now, we derive the remaining $n-1$ distinct roots of $X^n + 1 \bmod t$ as follows:

$(\omega^3)^n + 1 \equiv (\omega^{2n})\cdot \omega^n + 1 \equiv \omega^n + 1 \equiv t - 1 + 1 \equiv 0 \bmod t$

$(\omega^5)^n + 1 \equiv (\omega^{4n})\cdot \omega^n + 1 \equiv \omega^n + 1 \equiv t - 1 + 1 \equiv 0 \bmod t$

$\vdots$

$(\omega^{2n - 1})^n + 1 \equiv 0 \bmod t$ \\ \textcolor{red}{ $\rhd$ for any odd $k = 2j+1$, $(\omega^{2j+1})^n = (\omega^{2n})^j \cdot \omega^n = 1^j \cdot (-1) = -1$, so $(\omega^k)^n + 1 = 0$}

$ $

Note that $\omega, \omega^3, \omega^5, \cdots, \omega^{2n - 1}$ are all distinct values in $\mathbb{Z}_t^{\times}$, because $\textsf{Ord}_{\mathbb{Z}_t}(\omega) = 2n$. Thus, $\{\omega^{2i + 1}\}_{i=0}^{n-1}$ are $n$ distinct roots of $X^n + 1$. At the same time, since these are the roots of the cyclotomic polynomial $X^n + 1$, these are $n$ distinct primitive $(\mu=2n)$-th roots of unity.


$ $


We summarize our findings as follows:

\begin{tcolorbox}[title={\textbf{\tboxtheorem{\ref*{subsubsec:poly-vector-transformation-modulus}} Isomorphism between Polynomials and Vectors over Integers (Ring)}}]

\begin{itemize}
\item Suppose we have an $(n-1)$-degree polynomial ring $\mathbb{Z}_t[X] / F(X)$ where $F(X)$ is an $n$-degree polynomial having $n$ distinct roots $\{x_0, \cdots, x_{n-1}\}$, and an $n$-dimensional vector space $\mathbb{Z}_t^n$ (integers mod $t$) with vector $\vec{v}$.

$\sigma: f(X) \in \mathbb{Z}_t[X] / F(X) \text{ } \longrightarrow \text{ } (f(x_0), f(x_1), f(x_2), \cdots, f(x_{n-1})) \in \mathbb{Z}_t^n$

$ $

Then, $\sigma$ preserves isomorphism over the $(+, \cdot)$ operations:


$\sigma(f_a(X) + f_b(X)) = \sigma(f_a(X)) + \sigma(f_b(X))$

$\sigma(f_a(X) \cdot f_b(X)) = \sigma(f_a(X)) \odot \sigma(f_b(X))$

$ $

%\item If the $n$-degree polynomial $F(x)$ has a fewer number of roots than $n$, then isomorphism between the vectors and polynomials over the $(+, \cdot)$ operations holds without modulo $t$ as follows:

%$\sigma': f(X) \in \mathbb{Z}[X] / F(X) \text{ } \longrightarrow \text{ } (f(x_0), f(x_1), f(x_2), \cdots, f(x_{n-1})) \in \mathbb{Z}^n$

%$ $

\item Suppose we have the $(\mu=2n)$-th cyclotomic polynomial $X^n + 1 \bmod t$ such that $t$ is a prime and $t - 1$ is some multiple of $2n$, and $g$ is a generator of $\mathbb{Z}_t^{\times}$. Then, $n$ distinct roots of $X^n + 1$ (i.e., primitive $(\mu=2n)$-th roots of unity) can be efficiently computed as: $\{\omega^{2i + 1}\}_{i=0}^{n-1}$ where $\omega = g^{\frac{t-1}{2n}} \bmod t$. 


\end{itemize}
\end{tcolorbox}





\subsection{Isomorphism between Polynomials and Vectors over Complex Numbers}
\label{subsec:poly-vector-transformation-complex}

Here $n=2^r$ with $r\ge1$, so $n$ is even. All isomorphisms involving $\hat{\mathbb C}^n$ are real vector-space and real-algebra isomorphisms.


In Theorem~\ref*{subsec:polynomial-ring-basis} (\autoref{subsec:polynomial-ring-basis}), we learned the isomorphic mapping
\[
\sigma: f(X) \in \mathbb{Z}_t[X]/(X^n + 1) \longrightarrow (f(x_0),f(x_1),f(x_2), \cdots, f(x_{n-1})) \in \mathbb{Z}_t^n
\]
, where $x_0, x_1, x_2, \cdots, x_{n-1} \in \mathbb{Z}_t$ are the ($(\mu=2n)$-th primitive) roots of the cyclotomic polynomial $X^n + 1$, which are $\omega, \omega^3, \omega^5, \cdots, \omega^{2n-1}$, where $\omega$ can be any root of $X^n + 1$ (i.e., since each $\omega$ is a generator of all roots). In this subsection, we will demonstrate the isomorphism between a vector space and a polynomial ring over complex numbers as follows:

$\sigma_c: f(X) \in \mathbb{R}[X]/(X^n + 1) \longrightarrow (f(\omega),f(\omega^3),f(\omega^5), \cdots, f(\omega^{2n-1})) \in \mathbb{\hat{C}}^{n}$

$ $

, where $X \in \mathbb{C}$, and $\omega = e^{i\pi/n}$, the root (i.e., the primitive $(\mu=2n)$-th root) of the cyclotomic polynomial $X^n + 1$ over complex numbers (Theorem~\ref*{subsec:cyclotomic-theorem}.1 in \autoref{subsec:cyclotomic-theorem}). We define $\mathbb{\hat{C}}^{n}$ to be an $n$-dimensional \textit{special} vector space whose second-half elements of each vector are reverse-ordered conjugates of the first-half elements (e.g., $(v_0, v_1, \cdots, v_{\frac{n}{2}-1}, \overline{v}_{\frac{n}{2}-1}, \cdots, \overline{v}_1, \overline{v}_0 )$). 

\subsubsection{Isomorphism between $\mathbb{C}^{\frac{n}{2}}$ and $\hat{\mathbb{C}}^{n}$}
\label{subsec:poly-vector-transformation-complex-isomorphism1}

In this section, we treat both $\mathbb{C}^{\frac{n}{2}}$ and $\hat{\mathbb{C}}^{n}$ as vector spaces of complex vectors, where the mapping between them is defined as $\phi: \mathbb{C}^{\frac{n}{2}} \rightarrow \hat{\mathbb{C}}^{n}$ as follows:

$\phi(z_0, z_1, \dots, z_{\frac{n}{2}-1}) = (z_0, z_1, \dots, z_{\frac{n}{2}-1}, \overline{z_{\frac{n}{2}-1}}, \dots, \overline{z_1}, \overline{z_0})$


$ $

In other words, $\hat{\mathbb{C}}^{n}$ is a conjugate extension of $\mathbb{C}^{\frac{n}{2}}$ such that the first-half elements of $\hat{\mathbb{C}}^{n}$ are identical to those of $\mathbb{C}^{\frac{n}{2}}$, and the second-half elements of $\hat{\mathbb{C}}^{n}$ are reverse-ordered conjugates of $\mathbb{C}^{\frac{n}{2}}$. 

Now, we will demonstrate that the mapping $\phi$ between $\mathbb{C}^{\frac{n}{2}}$ and $\hat{\mathbb{C}}^{n}$ is an isomorphism.


$ $

\para{Bijective:} The mapping $\phi$ is trivially bijective. Every vector in $\mathbb{C}^{\frac{n}{2}}$ uniquely maps to a vector in $\hat{\mathbb{C}}^{n}$ by appending its reverse-ordered complex conjugates. Conversely, any vector in $\hat{\mathbb{C}}^{n}$ can uniquely map back to $\mathbb{C}^{\frac{n}{2}}$ by simply dropping the second half of its elements, establishing a perfect one-to-one correspondence.

\para{Homomorphic:} We will demonstrate that $\phi$ preserves homomorphism over addition $(+)$ and element-wise vector multiplication $(\odot)$. Given two vectors $\vec{u}, \vec{v} \in \mathbb{C}^{\frac{n}{2}}$:
\begin{itemize}
    \item \textbf{Addition:} $\phi(\vec{u} + \vec{v}) = \phi(\vec{u}) + \phi(\vec{v})$ holds because complex conjugation distributes over addition (i.e., $\overline{u + v} = \overline{u} + \overline{v}$).
    \item \textbf{Element-wise Multiplication:} $\phi(\vec{u} \odot \vec{v}) = \phi(\vec{u}) \odot \phi(\vec{v})$ holds because the conjugate of a product is the product of the conjugates (i.e., $\overline{u \cdot v} = \overline{u} \cdot \overline{v}$).
\end{itemize}

$ $

Therefore, the mapping $\phi: \mathbb{C}^{\frac{n}{2}} \rightarrow \hat{\mathbb{C}}^{n}$ is an isomorphism.

\begin{tcolorbox}[title={\textbf{Vector Space of Complex Vectors with either Real or Complex Scalars}}]

By mathematical definition, a vector space must satisfy \textit{closure under scalar multiplication}: for any vector $\vec{v}$ in the space and any scalar $c$ from the scalar field, the product $c \cdot \vec{v}$ must also lie in the same vector space.  Under this requirement, $\mathbb{C}^{\frac{n}{2}}$ is naturally a vector space of complex vectors with complex scalars (though the same is true with real scalars, since complex includes real). On the other hand, $\hat{\mathbb{C}}^{n}$ is a vector space of complex vectors \textit{strictly with real scalars}.

$ $

To see why, recall that a vector $\vec{v} = (v_0, v_1, \dots, v_{n-1}) \in \hat{\mathbb{C}}^n$ must satisfy the Hermitian-symmetry constraint: $v_{n-1-k} = \overline{v_k}$ for all $k$. Consider an example with $n = 4$:

$$\vec{v} = (1+i,\ 2+2i,\ 2-2i,\ 1-i) \in \hat{\mathbb{C}}^4$$

Now, let us multiply $\vec{v}$ by the complex scalar $c = i$:

$$i \cdot \vec{v} = (-1+i,\ -2+2i,\ 2+2i,\ 1+i)$$

For this result to still lie in $\hat{\mathbb{C}}^4$, the third entry must be the conjugate of the second entry. But $\overline{-2+2i} = -2-2i$, which is not equal to $2+2i$. So $i \cdot \vec{v} \notin \hat{\mathbb{C}}^4$, and closure fails.

$ $

More generally, for any nonzero $\vec{v} \in \hat{\mathbb{C}}^n$ and any scalar $c$, the product $c \cdot \vec{v}$ preserves Hermitian symmetry if and only if $\overline{c} = c$; that is, if and only if $c$ is a real number. Therefore, $\hat{\mathbb{C}}^n$ is a valid vector space only when restricted to real scalars, and this is why $\hat{\mathbb{C}}^n$ is a vector space of complex vectors strictly with real scalars.

$ $

Throughout this book, we will regard both $\mathbb{C}^{\frac{n}{2}}$ and $\hat{\mathbb{C}}^{n}$ as vector spaces of complex vectors with real scalars.

$ $

\textbf{Isomorphism of Vector Spaces:} For a mapping between vector spaces to be an isomorphism, it has to be not only bijective and homomorphic over the $(+, \cdot)$ operations, but also homomorphic over scalar multiplication. Therefore, we need to prove that $\phi(c \cdot \vec{u}) = c \cdot \phi(\vec{u})$. This is always true in our case because the scalar $c$ is a real number, meaning it is unaffected by complex conjugation (i.e., $\overline{c} = c$, and thus $\overline{c \cdot u_k} = c \cdot \overline{u_k}$). 

$ $

\textbf{Dimension of Vector Spaces:} The dimension of a vector space is defined by how many scalars it takes to fully describe it. Therefore, the dimension of $\mathbb{C}^{\frac{n}{2}}$ is $n$ because it has $\dfrac{n}{2}$ elements where each element requires 2 scalars to describe (one real and one imaginary part). The dimension of $\hat{\mathbb{C}}^{n}$ is also $n$ because the first $\dfrac{n}{2}$ elements requires $\dfrac{n}{2}\cdot 2$ scalars to describe, and the latter $\dfrac{n}{2}$ elements are simply reverse-ordered conjugates of the former $\dfrac{n}{2}$ elements.

\end{tcolorbox}


$ $



\subsubsection{Isomorphism between $\hat{\mathbb{C}}^{n}$ and $\mathbb{R}[X] / (X^n + 1)$}
\label{subsec:poly-vector-transformation-complex-isomorphism2}

Now, we will demonstrate that $\sigma_c$ is an isomorphism (i.e., bijective and homomorphic) between $\hat{\mathbb{C}}^{n}$ and $\mathbb{R}[X] / (X^n + 1)$ by applying the same reasoning as described in the beginning of \autoref{subsec:poly-vector-transformation}. 

\begin{figure}[h!]
    \centering
  \includegraphics[width=0.3\linewidth]{figures/cyclotomic-polynomial.pdf}
  \caption{An illustration of the four roots of the 8th cyclotomic polynomial $x^4 + 1$}
  \label{fig:cyclotomic-polynomial}
\end{figure}

$ $

\para{Bijective:} Based on Euler's formula $e^{i\theta} = \cos\theta + i\cdot\sin\theta$ (\autoref{subsec:euler}), we can derive the following arithmetic relations: $\omega = \overline{\omega^{2n-1}}, \text{ } \omega^3 = \overline{\omega^{2n-3}}, \cdots, \omega^{n-1} = \overline{\omega^{n+1}}$. In other words, the one-half roots are conjugates of the other-half roots. This can also be pictorially understood based on a complex plane in \autoref{fig:cyclotomic-polynomial}, where red arrows represent the roots of the 8th cyclotomic polynomial $X^4 + 1$, comprising imaginary number and real number components. As shown in this figure, one half of the red arrows (i.e., roots) are a reflection of the other half on the $x$-axis (i.e., real number axis). This means that we can express these roots as an $n$-dimensional vector whose elements are the roots of $X^n + 1$, such that its second-half elements are a reverse-ordered conjugate of the first-half elements. Based on this vector design, the $\sigma_c$ mapping can be re-written as follows:

$\sigma_c(f(X)) = (f(\omega),f(\omega^3),f(\omega^5), \cdots, f(\omega^{n-1}), f(\overline{\omega^{n-1}}), f(\overline{\omega^{n-3}}), \cdots, f(\overline{\omega^3}), f(\overline{\omega}))$

$ $

Since $f(\overline{X}) = \overline{f(X)}$ (because $f(X)$ has strictly real coefficients), we can rewrite $\sigma_c$ as: 

$\sigma_c(f(X)) = (f(\omega),f(\omega^3),f(\omega^5), \cdots, f(\omega^{n-3}), f(\omega^{n-1}), \overline{f(\omega^{n-1})}, \overline{f(\omega^{n-3})}, \cdots, \overline{f(\omega^3)}, \overline{f(\omega)})$

$ $

This structure of vector exactly aligns with the definition of $\hat{\mathbb{C}}^n$: the second half of the elements of the $n$-dimensional vector $\vec{\hat v}$ is a reverse-ordered conjugate of the first half.

To prove bijectivity, first note that distinct real polynomials of degree less than $n$ cannot agree at the $n$ distinct roots, by Lagrange interpolation (Theorem~\ref*{sec:polynomial-interpolation} in \autoref{sec:polynomial-interpolation}). For surjectivity, take any $\vec v\in\hat{\mathbb C}^n$ and interpolate its values at the ordered roots to obtain a unique $f\in\mathbb C[X]$ of degree less than $n$. Conjugating the coefficients of $f$ produces another polynomial with the same evaluations, because both nodes and values occur in reversed conjugate pairs. Uniqueness forces that conjugated polynomial to equal $f$, so every coefficient of $f$ is real. Hence every $\vec v\in\hat{\mathbb C}^n$ has a real-polynomial preimage.

$ $

\para{Homomorphic:} $\sigma_c$ is homomorphic, because based on the reasoning shown in \autoref{subsec:poly-vector-transformation}, the relations $\sigma_c(f_a(X) + f_b(X)) = \sigma_c(f_a(X)) + \sigma_c(f_b(X))$ and $\sigma_c(f_a(X) \cdot f_b(X)) = \sigma_c(f_a(X)) \odot \sigma_c(f_b(X))$ mathematically hold regardless of whether the type of $X$ is a modulo integer or a complex number. Additionally, for any real scalar $c \in \mathbb{R}$, $\sigma_c(c \cdot f(X)) = c \cdot \sigma_c(f(X))$ holds because scaling the polynomial and then evaluating it is the same as evaluating it and then scaling the result.

$ $

Since $\sigma_c$ is both bijective and homomorphic over the $(+, \cdot)$ operations, it is an isomorphism. 


\begin{tcolorbox}[title={\textbf{\tboxtheorem{\ref*{subsec:poly-vector-transformation-complex}} Isomorphism between Polynomials and Vectors over Complex Numbers}}]

The following mapping $\sigma_c$ between polynomials and vectors over complex numbers is an isomorphism:

$\sigma_c: f(X) \in \mathbb{R}[X]/(X^n + 1) \longrightarrow (f(\omega),f(\omega^3),f(\omega^5), \cdots, f(\omega^{2n-1})) \in \hat{\mathbb{C}}^{n} \text{ } (\cong \mathbb{C}^{\frac{n}{2}})$

$ $

, where $\omega = e^{i\pi/n}$, the root (i.e., the primitive $(\mu=2n)$-th root) of the cyclotomic polynomial $X^n + 1$ over complex numbers, and $\hat{\mathbb{C}}^{n}$ is an $n$-dimensional vector space of complex vectors (with real scalars) whose second-half elements are reverse-ordered conjugates of the first-half elements. 

\end{tcolorbox}





\subsection{Transforming Basis between Polynomial Ring and Vector Space of Complex Vectors with Real Scalars}
\label{subsec:poly-vector-basis-transfer}


A linear isomorphism over a fixed scalar field sends a basis to a basis: linearity preserves all linear combinations, injectivity preserves linear independence, and surjectivity preserves spanning. The evaluation map $\sigma$ above is linear over $\mathbb Z_t$ with prime $t$, and $\sigma_c$ is linear over $\mathbb R$. Those are the scalar fields used in the following basis-transfer argument.

The same principle holds between a polynomial ring and vector space of complex numbers (with real scalars). Given the polynomial ring $\mathbb{R}[X]/(x^n + 1)$, the most intuitive way to set up a basis of $\mathbb{R}[X]/(x^n + 1)$ is as follows:

$f_0(X) = 1$

$f_1(X) = X$

$f_2(X) = X^2$

$\vdots$

$f_{n-1}(X) = X^{n-1}$

$ $

These $n$ polynomials are linearly independent, because each polynomial exclusively has its own unique exponent term, whereas one term cannot be expressed by a linear combination of the other terms. Also, these $n$ polynomials span the polynomial ring $\mathbb{R}[X]/(x^n + 1)$, because each polynomial's scalar multiplication can express any coefficient value of its own exponent term, and summing all such polynomials can express any polynomial in the polynomial ring $\mathbb{R}[X]/(X^n + 1)$.


Now, we will apply the $\sigma_c$ mapping to the above $n$ polynomials that are a basis of the $(n-1)$-degree polynomial ring $\mathbb{R}[X]/(x^n + 1)$. Then, according to the principle of polynomial-to-vector basis transfer (explained in Theorem~\ref*{subsec:polynomial-ring-basis} in \autoref{subsec:polynomial-ring-basis}), we can use these $n$ polynomials (i.e., the basis of the $(n-1)$-degree polynomial ring) and the isomorphic polynomial-to-vector mapping $\sigma_c$ to compute the basis of the $n$-dimensional special vector space $\hat{\mathbb C}^{n}$ as follows: 

$W = \begin{bmatrix}
\sigma_c(f_0(X)) \\
\sigma_c(f_1(X)) \\
\sigma_c(f_2(X)) \\
\vdots \\
\sigma_c(f_{n-1}(X)) \\
\end{bmatrix}
= \begin{bmatrix}
\sigma_c(1) \\
\sigma_c(X) \\
\sigma_c(X^2) \\
\vdots \\
\sigma_c(X^{n-1}) \\
\end{bmatrix}$
$= \begin{bmatrix}
1 & 1 & \cdots & 1 & 1\\
(\omega) & (\omega^3) & (\omega^5) & \cdots & (\omega^{2n-1})\\
(\omega)^2 & (\omega^3)^2 & (\omega^5)^2 & \cdots & (\omega^{2n-1})^2\\
\vdots & \vdots & \vdots & \ddots & \vdots \\
(\omega)^{n-1} & (\omega^3)^{n-1} & (\omega^5)^{n-1} & \cdots & (\omega^{2n-1})^{n-1}\\
\end{bmatrix}$

\text{ } \text{ }  $= \begin{bmatrix}
1 & 1 & 1 & \cdots & 1\\
(\omega) & (\omega^3) & \cdots & (\overline{\omega})^3 & (\overline{\omega})\\
(\omega)^2 & (\omega^3)^2 & \cdots & (\overline{\omega}^3)^2 & (\overline{\omega})^2\\
\vdots & \vdots & \vdots & \ddots & \vdots \\
(\omega)^{n-1} & (\omega^3)^{n-1} & \cdots & (\overline{\omega}^3)^{n-1} & (\overline{\omega})^{n-1}\\
\end{bmatrix}$

$ $

$W$ is a valid basis of the $n$-dimensional special vector space $\hat{\mathbb C}^{n}$. 



\begin{tcolorbox}[title={\textbf{\tboxtheorem{\ref*{subsec:poly-vector-transformation-complex}} Transforming Basis between Polynomial Ring and Vector Space}}]

If $n$ polynomials form a basis of an $(n-1)$-degree polynomial ring and they are converted into $n$ distinct vectors via an isomorphic mapping $\sigma$ (or $\sigma_c$ in the case of the complex number domain) from the $(n-1)$-degree polynomial ring to the $n$-dimensional vector space, then those converted $n$ vectors form a basis of the $n$-dimensional vector space. 

\end{tcolorbox}
